% !TeX root = ../MAgent_Theory_2.0_supplement.tex
\chapter{Algebra of Selective Exponential Links}
\label{supp:geometry}

This chapter gives supporting matrix identities for interactions written in a reference vector space. All spaces are finite dimensional and real. The results concern arbitrary finite matrices and do not require Gaussian laws, diagonalizability, or invertibility of the generator. A statistical interpretation enters only when the matrices are assigned to interfaces or made functions of law parameters.

\section{The low-rank exponential}

Let $O:\R^r\to\R^d$ be a receiving map and $T:\R^d\to\R^r$ a sending map. Define the entire matrix function
\begin{equation}
\varphi_1(Z)=\sum_{n=1}^{\infty}\frac{Z^{n-1}}{n!}.
\end{equation}
When $Z$ is invertible this equals $(e^Z-I)Z^{-1}$, but the series defines it also at singular $Z$.

\lemmaheading{Exponential through the message space}{lem:sgeom-lowrank}
The following identity holds for all $O,T$:
\begin{equation}
e^{OT}=I_d+O\varphi_1(TO)T.
\label{eq:sgeom-lowrank}
\end{equation}
\begin{proof}
Associativity gives $(OT)^n=O(TO)^{n-1}T$ for every integer $n\geq1$. Absolute convergence of the finite-dimensional exponential series permits summation and yields \cref{eq:sgeom-lowrank}.
\end{proof}
In particular, the correction to identity has image in $\operatorname{im}O$ and vanishes on $\ker T$. The image of the full exponential need not be low dimensional. Indeed its inverse is $e^{-OT}$, so the identity contribution preserves all input dimensions even when the communicated correction is selective.

The identity also applies to several interfaces simultaneously. For maps $O_h,T_h$, stack $O=[O_1\ \cdots\ O_H]$ and stack the $T_h$ vertically to form $T$. Then $OT=\sum_hO_hT_h$ and the same formula holds. The smaller matrix $TO$ includes all cross-interface products $T_hO_k$; replacing it by its diagonal blocks would require additional vanishing conditions.

\propositionheading{Rank-one closed form and determinant}{prop:sgeom-rankone}
For column vectors $u,v\in\R^d$, put $L=uv^T$ and $t=v^Tu$. Then
\begin{equation}
e^L=I_d+\varphi_1(t)L,\qquad
\det e^L=e^t.
\end{equation}
If $t=0$, then $L^2=0$ and $e^L=I_d+L$.
\begin{proof}
The first identity is \cref{eq:sgeom-lowrank} with a one-dimensional message space. For any square matrix $A$, Jacobi's determinant identity applied to $X(s)=e^{sA}$ gives
$\frac{d}{ds}\det X(s)=\Tr(A)\det X(s)$.
Since $X(0)=I$, its determinant is $e^{s\Tr(A)}$. Set $A=L$, $s=1$, and use $\Tr(uv^T)=v^Tu$. Finally $L^2=u(v^Tu)v^T$ proves the nilpotent case.
\end{proof}
This result distinguishes a singular generator from a singular interaction: a rank-one generator gives an invertible exponential in every dimension. It does not assert that every invertible real matrix is one exponential.

\section{Differentiating a law-dependent link}

\lemmaheading{Derivative of the matrix exponential}{lem:sgeom-frechet}
For square matrices $L,H$, the Fr\'echet derivative is
\begin{equation}
D\exp_L[H]=\int_0^1e^{(1-t)L}He^{tL}dt.
\label{eq:sgeom-frechet}
\end{equation}
\begin{proof}
Let $X_\epsilon(t)=e^{t(L+\epsilon H)}$. Differentiating its matrix differential equation with respect to $\epsilon$ at zero gives
$\dot Y=LY+He^{tL}$, $Y(0)=0$.
Multiplying by $e^{-tL}$ and integrating gives
$Y(t)=\int_0^t e^{(t-s)L}He^{sL}ds$.
Evaluation at $t=1$ proves the formula. Smooth dependence is also obtained directly by differentiating the locally uniformly convergent exponential series.
\end{proof}
If $[L,H]=0$, this derivative reduces to $e^LH$. Without that hypothesis the replacement is invalid in general. For a link $\Theta=U_i e^L U_j^{-1}$ with fixed endpoint frames, its variation is $U_iD\exp_L[dL]U_j^{-1}$. If the frames vary, the product rule supplies their two additional terms, including $d(U_j^{-1})=-U_j^{-1}(dU_j)U_j^{-1}$.

For the moment-dependent family
$L=\lambda\mu_i\mu_j^TM_j^{-1}$, with fixed scalar $\lambda$, ordinary differentiation gives
\begin{equation}
\begin{split}
dL=\lambda\bigl[&(d\mu_i)\mu_j^TM_j^{-1}
+\mu_i(d\mu_j)^TM_j^{-1}\\
&-\mu_i\mu_j^TM_j^{-1}(dM_j)M_j^{-1}\bigr].
\end{split}
\label{eq:sgeom-moment-derivative}
\end{equation}
Here a differentiable family of laws is assumed to have differentiable first and second moments; differentiation under their defining integrals requires the corresponding domination or integrability hypotheses. Positive-definite covariance makes inversion smooth. Combining \cref{eq:sgeom-moment-derivative,eq:sgeom-frechet} gives the complete link derivative on this admitted family, with no assumption about the shape of its densities.

The resulting terms enter a variational derivative whenever an interaction energy reads the link. Holding the generator fixed suppresses them and defines a different update convention. This algebra separates the constitutive choice of a law-dependent interaction from the kinetic choice used to evolve the laws; neither choice follows from the exponential parameterization alone.
